\documentclass[a4paper,12pt]{report}
\usepackage[utf8]{inputenc}
\usepackage{graphicx}
\usepackage{float}
\usepackage{tabularx}
\usepackage{makecell}
\usepackage{titlesec}
\usepackage{fancyhdr}
\usepackage{lastpage}
\usepackage{xurl}
\usepackage{hyperref}
\usepackage{geometry}
\usepackage{color}
\usepackage{microtype}
\usepackage{enumerate}
\usepackage{listings}
\usepackage[table,dvipsnames]{xcolor}
\setcounter{tocdepth}{5}
\setcounter{secnumdepth}{5} %SONO I DUE COMANDI PER AVERE LE SUBSUBSECTION NUMERATE E PRESENTI NELL'INDICE
\renewcommand{\contentsname}{Indice} %QUESTO SERVE PER AVERE L'INDICE CON IL NOME CHE VOGLIO IO
\titleformat{\chapter}[display]
{\normalfont\bfseries}{}{0pt}{\LARGE} %QUESTO SERVE PER AVERE SOLO IL NOME DEL CAPITOLO CHE VOGLIO IO
\titlespacing*{\chapter}{0cm}{0cm}{0.2cm}
\geometry{
	left=20mm,
	right=20mm,
}
\fancypagestyle{plain}{
	\fancyhf{}
	\lhead{\includegraphics[width=3cm]{../immagini/minilogo.jpg}}
	\chead{}
	\rhead{\fontsize{12}{10}Linee Guida}
	\lfoot{}
	\cfoot{\thepage\ di \pageref*{LastPage}}
	\rfoot{}
}

\definecolor{atomlightorange}{rgb}{0.88,0.76,0.55}
\definecolor{atomdarkgrey}{RGB}{59,62,75}

\usepackage{tikz}

% set listings
\lstset{%
    basicstyle=\footnotesize\ttfamily\color{atomlightorange},
    framesep=20pt,
		belowskip=10pt,
		aboveskip=10pt
}

% add frame environment
\usepackage[%
    framemethod=tikz,
    skipbelow=8pt,
    skipabove=13pt
]{mdframed}
\mdfsetup{%
    leftmargin=0pt,
    rightmargin=0pt,
    backgroundcolor=atomdarkgrey,
    middlelinecolor=atomdarkgrey,
    roundcorner=6
}

\usepackage{etoolbox}% >= v2.1 2011-01-03
\BeforeBeginEnvironment{lstlisting}{\begin{mdframed}\vspace{-0.7em}}
\AfterEndEnvironment{lstlisting}{\vspace{-0.5em}\end{mdframed}}

% needed for \lstcapt
\def\ifempty#1{\def\temparg{#1}\ifx\temparg\empty}

% make new caption command for listings
\usepackage{caption}
\newcommand{\lstcapt}[2][]{%
    \ifempty{#1}%
        \captionof{lstlisting}{#2}%
    \else%
        \captionof{lstlisting}[#1]{#2}%
    \fi%
    \vspace{0.75\baselineskip}%
}

\hypersetup{
    colorlinks=true,
    linkcolor=black,
    filecolor=black,
    urlcolor=blue,
		citecolor=black,
}

\pagestyle{plain}

\makeatletter
	 \def\thebibliography#1{\chapter*{Bibliografia\@mkboth
		 {Bibliografia}{Bibliografia}}\list
		 {[\arabic{enumi}]}{\settowidth\labelwidth{[#1]}\leftmargin\labelwidth
 \advance\leftmargin\labelsep
 \usecounter{enumi}}
 \def\newblock{\hskip .11em plus .33em minus .07em}
 \sloppy\clubpenalty4000\widowpenalty4000
 \sfcode`\.=1000\relax}
	 \makeatother


\begin{document}

\makeatletter
\begin{titlepage}
    \begin{center}
    \vspace*{-4,0cm}
    \author{Jawa Druids}
    \title{Linee Guida}
    \date{} %LASCIARE QUESTO CAMPO VUOTO, SE LO TOLGO STAMPA LA DATA CORRENTE
    \includegraphics[width=0.7\linewidth]{../immagini/DRUIDSLOGO.jpg}\\[4ex]
    {\huge \bfseries  \@title }\\[2ex]
    {\LARGE  \@author}\\[50ex]
    \vspace*{-8,0cm}
    \begin{table}[H]
        \centering
        \begin{tabular}{c|c}
            \textbf{Versione} & x.x.x \\
            \textbf{Data approvazione} & xx-xx-xxxx\\
            \textbf{Responsabile} & Nome Cognome\\
            \textbf{Redattori} & Andrea Dorigo \\
            \textbf{Verificatori} & Nome Cognome \\
        %MAKECELL SERVE PER POI ANDARE A CAPO ALL'INTERNO DELLA CELLA
            \textbf{Stato} & Redazione in corso\\
            \textbf{Lista distribuzione} & \makecell{Jawa Druids \\ Tullio \\ NomeAziendaProponente}\\
            \textbf{Uso} & Interno
        \end{tabular}
    \end{table}
					\vspace{.8cm}
	\hfill \break
    \fontsize{17}{10}\textbf{Sommario}\\
		\vspace{.3cm}
    Linee guide da seguire per una sessione lavorativa standard con git flow, Trello e Google Sheet.
    \end{center}
\end{titlepage}
\makeatother

\begin{comment}
	\quad
	\begin{center}
	\LARGE\textbf{Registro delle modifiche}
	\end{center}
	\def\tabularxcolumn#1{m{#1}}
	{\rowcolors{2}{RawSienna!90!RawSienna!20}{RawSienna!70!RawSienna!40}

	\begin{center}
	\renewcommand{\arraystretch}{1.4}
	\begin{longtable}[c]{|p{5.5cm}|p{3cm}|p{2.1cm}|p{2cm}|p{2cm}|}
	\hline
	\rowcolor{airforceblue}
	\textbf{Modifica} & \textbf{Autore} & \textbf{Ruolo} & \textbf{Data} & \textbf{Versione}\\
	\hline
	\textit{Approvazione del documento per RR.} & Andrea Cecchin & \textit{Responsabile} & 10-01-2021 & v1.0.0 \\
	\hline
	\textit{Verifica intero documento.} & Alfredo Graziano & \textit{Verificatore} & 09-01-2021 & v0.2.0 \\
	\hline
	\textit{Aggiunta  \S~\ref{requisitiConsiderazioni}.} & Igli Mezini & \textit{Analista} & 07-01-2021 & v0.1.5 \\
	\hline
	\textit{Aggiunte sezioni \S~\ref{RequisitiTracciamentoDeiRequisitiFonte}, \S~\ref{RequisitiTracciamentoDeiRequisitiFonteRequisito}.} & Igli Mezini & \textit{Analista} & 06-01-2021 & v0.1.4 \\
	\hline
	\textit{Aggiunte sezioni \S~\ref{descrizioneGeneraleCaratteristicheUtente}, \S~\ref{RequisitiPrestazionali}, \S~\ref{RequisitiDiQualita}, \S~\ref{RequisitiVincolo} .} & Igli Mezini & \textit{Analista} & 05-01-2021 & v0.1.3 \\
	\hline
	\textit{Modifica  \S~\ref{RequisitiFunzionali} .} & Andrea Cecchin & \textit{Analista} & 05-01-2021 & v0.1.2 \\
	\hline
	\textit{Inizio stesura \S~\ref{Requisiti}.} & Andrea Cecchin & \textit{Analista} & 04-01-2021 & v0.1.1 \\
	\hline
	\textit{Verifica \S~\ref{Introduzione},\S~\ref{descrizioneGenerale}.} & Mattia Cocco & \textit{Verificatore} & 28-12-2020 & v0.1.0 \\
	\hline
	\textit{Stesura \S~\ref{fasiProgetto}.} & Andrea Cecchin & \textit{Analista} & 28-12-2020 & v0.0.4 \\
	\hline
	\textit{Aggiunte sezioni \S~\ref{fasiProgettoAquisizioneDati}, \S~\ref{fasiProgettoElaborazioneDati}.} & Andrea Cecchin & \textit{Analista} & 10-12-2020 & v0.0.3 \\
	\hline
	\textit{Aggiunte sezioni \S~\ref{descrizioneGeneraleCaratteristicheProdotto}, \S~\ref{descrizioneGeneraleFunzionalitàGenerali}.} & Andrea Cecchin & \textit{Analista} & 10-12-2020 & v0.0.2 \\
	\hline
	\textit{Inizio stesura documento \S~\ref{Introduzione}.} & Andrea Cecchin & \textit{Analista} & 07-12-2020 & v0.0.1 \\
	\hline
	\end{longtable}
	\end{center}
\end{comment}

\quad
\begin{center}
	\LARGE\textbf{Registro delle modifiche}
\end{center}
\def\tabularxcolumn#1{m{#1}}
{\rowcolors{2}{RawSienna!90!RawSienna!20}{RawSienna!70!RawSienna!40}

\begin{center}
	\renewcommand{\arraystretch}{1.4}
	\begin{longtable}[c]{|p{2cm-1\tabcolsep}|p{2cm}|p{3cm-2\tabcolsep}|p{2,5cm-2\tabcolsep}|p{3,5cm-2\tabcolsep}|p{3cm}|}
		\hline
		\rowcolor{airforceblue}
		\makecell[c]{\textbf{Versione}} & \makecell[c]{\textbf{Data}} & \makecell[c]{\textbf{Autore}} & \makecell[c]{\textbf{Ruolo}} & \makecell[c]{\textbf{Modifica}} & \makecell[c]{\textbf{Verificatore}} \\
		\hline
		\centering v4.0.0 & 2021-04-29 & Andrea Cecchin & \centering \textit{Responsabile} & \textit{Approvazione del documento per RQ} & \makecell[c]{-}\\
		\hline
		\centering v3.3.0 & 2021-04-28 & \centering - & \centering - & \textit{Revisione complessiva del documento} & \makecell[c]{Emma Roveroni}\\
		\hline
		\centering v3.2.1 & 2021-04-27 & Andrea Cecchin & \centering \textit{Analista} & \textit{Revisione dei casi d'uso} & \makecell[c]{Andrea Dorigo}\\
		\hline
		\centering v3.2.0 & 2021-04-18 & \centering - & \centering \textit{-} & \textit{Revisione complessiva a seguito dei cambiamenti} & \makecell[c]{Emma Roveroni}\\
		\hline
		\centering v3.1.2 & 2021-04-18 & Igli Mezini & \centering \textit{Analista} & \textit{Sistemato capitolo \ref{Requisiti}} & \makecell[c]{Alfredo Graziano}\\
		\hline
		\centering v3.1.1 & 2021-04-17 & Andrea Cecchin & \centering \textit{Analista} & \textit{Correzzione complessiva del capitolo \ref{CasiDUso}} & \makecell[c]{Alfredo Graziano}\\
		\hline
		\centering v3.1.0 & 2021-04-09 & \centering - & \centering \textit{-} & \textit{Revisione complessiva a seguito dei cambiamenti} & \makecell[c]{Emma Roveroni}\\
		\hline
		\centering v3.0.2 & 2021-04-08 & Mattia Cocco & \centering \textit{Analista} & \textit{Correzione descrizione/scenario di \S~\ref{CasiDUsoCasiDUsoTraUnUtenteEIlFrontEndElencoCasiDUsoUC3ZoomDellaHeatMap}, \S~\ref{CasiDUsoCasiDUsoTraUnUtenteEIlFrontEndElencoCasiDUsoUC311SpostamentoDelCentroDellaMappa}, \S~\ref{CasiDUsoCasiDUsoTraUnUtenteEIlFrontEndElencoCasiDUsoUC312VisualizzazioneDelPopupDiUnPuntoDiInteresse}, \S~\ref{CasiDUsoCasiDUsoFacoltativiTraUnUtenteEIlFrontEndElencoCasiDUsoUC15NotificaViaEmailDiUnaCittaSelezionata}} &  \makecell[c]{Emma Roveroni} \\
		\hline
		\centering v3.0.1 & 2021-04-06 & Mattia Cocco & \centering \textit{Analista} & \textit{Correzione diagrammi casi d'uso UC3-8-15} & \makecell[c]{Emma Roveroni}\\
		\hline
		\centering v3.0.0 & 2021-03-15 & Margherita Mitillo & \centering \textit{Responsabile} & \textit{Approvazione del documento per RP} & \makecell[c]{-}\\
		\hline
		\centering v2.2.0 & 2021-03-09 &\centering - & \centering \textit{-} &\textit{Revisione complessiva a seguito dei cambiamenti} & Mattia Cocco\\
		\hline
		\centering v2.1.1 & 2021-03-07 & Margherita Mitillo & \centering{ \textit{Analista} }&\textit{Rielaborazione ed aggiornamento del capitolo \S~\ref{Requisiti}} &  Alfredo Graziano  \\
		\hline
		\centering v2.1.0 & 2021-03-05 &\centering - & \centering \textit{-} &  \textit{Revisione complessiva a seguito dei cambiamenti}  & Andrea Dorigo\\
		\hline
		\centering v2.0.1 & 2021-03-04 & Margherita Mitillo & \centering \textit{Analista} &\textit{Rielaborazione ed aggiornamento del capitolo \S~\ref{CasiDUso}} & Alfredo Graziano  \\
		\hline
		\centering v2.0.0 & 2021-03-01 & Andrea Cecchin & \centering \textit{Responsabile} &\textit{Approvazione del documento per la riconsegna della RR} &  \makecell[c]{-}\\
		\hline
		\centering v1.2.0 & 2021-02-28 & Andrea Dorigo & \centering \textit{Analista} &\textit{Rielaborazione del capitolo \S~\ref{DescrizioneGenerale}} & Andrea Cecchin  \\
		\hline
		\centering v1.1.0 & 2021-02-15 & \centering - & \centering - &\textit{Verifica del documento} & Andrea Dorigo  \\
		\hline
		\centering v1.0.2 & 2021-02-12 & Andrea Cecchin & \centering \textit{Analista} &\textit{Correzioni e aggiunte al capitolo \S~\ref{Requisiti}} & Mattia Cocco  \\
		\hline
		\centering v1.0.1 & 2021-02-11 & Andrea Cecchin & \centering \textit{Analista} &\textit{Ristesura capitolo \S~\ref{CasiDUso}} & Mattia Cocco  \\
		\hline
		\centering v1.0.0 & 2021-01-10 & Andrea Cecchin & \centering \textit{Responsabile} &\textit{Approvazione del documento per RR.} &  \makecell[c]{-}  \\
		\hline
		\centering v0.2.0 & 2021-01-09 & \centering - & \centering - &\textit{Verifica intero documento.} & Alfredo Graziano  \\
		\hline
		\centering v0.1.5 & 2021-01-07 & Igli Mezini & \centering \textit{Analista} & \textit{Aggiunta  \S~\ref{RequisitiConsiderazioni}.} & \makecell[c]{-}  \\
		\hline
		\centering v0.1.4 & 2021-01-06 & Igli Mezini & \centering \textit{Analista} & \textit{Aggiunte \S~\ref{RequisitiTracciamentoDeiRequisitiFonte}, \S~\ref{RequisitiTracciamentoDeiRequisitiFonteRequisito}.} &   \makecell[c]{-}  \\
		\hline
		\centering v0.1.3 & 2021-01-05 & Igli Mezini & \centering \textit{Analista} & \textit{Aggiunte \S~\ref{DescrizioneGeneraleCaratteristicheUtente}, \S~\ref{RequisitiPrestazionali}, \S~\ref{RequisitiDiQualita}, \S~\ref{RequisitiDiVincolo} .} &  \makecell[c]{-}  \\
		\hline
		\centering v0.1.2 & 2021-01-05 & Andrea Cecchin & \centering \textit{Analista} & \textit{Modifica  \S~\ref{RequisitiFunzionali} .} &    \makecell[c]{-} \\
		\hline
		\centering v0.1.1 & 2021-01-04 & Andrea Cecchin & \centering \textit{Analista} & \textit{Inizio stesura \S~\ref{Requisiti}.} & \makecell[c]{-}\\
		\hline
		\centering v0.1.0 & 2020-12-28 & \centering - & \centering - & \textit{Verifica \S~\ref{Introduzione},\S~\ref{DescrizioneGenerale}.} &  Mattia Cocco\\
		\hline
		\centering v0.0.4 & 2020-12-28 & Andrea Cecchin & \centering \textit{Analista} & \textit{Fine stesura capitolo \S~3} & \makecell[c]{-} \\
		\hline
		\centering v0.0.3 & 2020-12-10 & Andrea Cecchin & \centering \textit{Analista} & \textit{Aggiunte \S~3.1 e \S~3.2} & \makecell[c]{-} \\
		\hline
		\centering v0.0.2 & 2020-12-10 & Andrea Cecchin & \centering \textit{Analista} & \textit{Aggiunte \S~\ref{DescrizioneGeneraleCaratteristicheProdotto}, \S~\ref{DescrizioneGeneraleFunzionalitàDelProdotto}} &  \makecell[c]{-} \\
		\hline
		\centering v0.0.1 & 2020-12-07 & Andrea Cecchin & \centering \textit{Analista} & \textit{Inizio stesura documento \S~\ref{Introduzione}} &  \makecell[c]{-} \\
		\hline

	\end{longtable}
\end{center}

\tableofcontents{}
\chapter{Introduzione}

Lo scopo di questo documento è la stesura di un elenco di linee guida ed esempi che i componenti sono incoraggiati a seguire per migliorare l'efficacia della collaborazione.
A differenza delle Norme di Progetto, questo documento è redatto in un linguaggio più informale per facilitarne la comprensione, con spezzati di codice a scopo esemplificativo e riferimenti esterni sulle best practices da seguire.\\
Il documento può essere soggetto a modifiche ed aggiunte per tutta la durata del progetto.

% \lstset{style=mystyle}


\chapter{Git}

\section{Aggiunta o modifica di files}
Poniamo che un componente del gruppo voglia aggiungere o modificare uno o più files nella repository del progetto.
Il primo passo è creare una nuovo feature branch tramite il comando
\begin{lstlisting}
git flow feature start analisi-uber
\end{lstlisting}
Questo crea un branch locale a partire dal develop chiamato \texttt{feature/analisi-uber} e vi esegue il checkout automaticamente; tutti i prossimi commit e push avverranno su questo branch, evitando l'immissione di errori o file ancora incompleti o non verificati sul \texttt{develop}.
\textbf{Per rendere possibile la collaborazione sullo stesso branch} è possibile pubblicarlo in remoto con il comando
\begin{lstlisting}
git flow feature publish analisi-uber
\end{lstlisting}
Chiunque desideri lavorare su questo branch necessita di immettere
\begin{lstlisting}
 git checkout -b feature/analisi-uber origin/feature/analisi-uber
\end{lstlisting}
che crea un branch locale dallo stesso nome e lo sincronizza con quello remoto.
Al completamento \textbf{e verifica} della nuova funzionalità, è possibile richiudere il branch nel \texttt{develop} ed eliminarlo tramite il comando
\begin{lstlisting}
git flow feature finish analisi-uber
\end{lstlisting}
Verrà automaticamente generato un \texttt{merge commit}.
Il responsabile di progetto si riserva la pubblicazione delle \texttt{release}.

\section{Link utili e best practices sull'utilizzo di git}
Potete trovare un elenco dei comandi di git flow più importanti e relative funzioni
\href{https://danielkummer.github.io/git-flow-cheatsheet/}{qui} \cite{GitFlowCheatseet}. \\
Un buon riassunto sull'utilizzo di git flow si può trovare \href{https://www.atlassian.com/git/tutorials/comparing-workflows/gitflow-workflow}{a questo link} \cite{GitFlowWorkflowTutorial}.\\
Dell'altro materiale riguardante le best practice da seguire nei commit \href{https://www.theserverside.com/video/Follow-these-git-commit-message-guidelines}{a questo indirizzo} \cite{GitCommitGuidelines}. \\
Fra le molte pratiche di cui tener conto nella scrittura del titolo del commit si evidenzia:
\begin{itemize}
\item esso deve contenere la descrizionee di cosa è stato fatto e perchè, ma non come.
\item l'utilizzo dell'imperativo;
\end{itemize}
Per quest'ultimo punto vi riporto un paragrafo del link qui sopra che ne facilita la comprensione (è presente un errore nel testo dell'articolo, che ho corretto qui sotto):
\begin{lstlisting}
Fixed the fencepost error //bad

Fixing the fencepost error //bad

Fix the fencepost error //good
\end{lstlisting}
\textit{The imperative mood is the one git commit message guideline that developers tend to violate most often. A good rule of thumb is that a git commit message can be appended to the statement "If applied, this commit will …." The resulting sentence should make grammatical sense. As you can see from the following three examples, gerunds and past tense commits fail the test, while the imperative tense does not:}
\begin{lstlisting}
If applied, this commit will Fixed the fencepost error //bad

If applied, this commit will Fixing the fencepost error //bad

If applied, this commit will Fix the fencepost error //good
\end{lstlisting}

\chapter{Trello}

Trello è la piattaforma scelta dal gruppo per la gestione del progetto, data la famigliarit\`{a} di certi componenti del gruppo con esso e la comprovata efficacia.
Qui sono raccolti tutti gli stadi futuri/presenti/passati del progetto, e l'attivit\`{a} che ogni componente sta svolgendo.

La \textbf{creazione} e lo \textbf{spostamento} delle schede è riservata all'amministratore o al responsabile (o quantomeno è richiesta la loro approvazione).
L'\textbf{assegnazione} delle schede è riservata al responsabile (o quantomeno è richiesta la sua approvazione).

\chapter{Sessione lavorativa standard}

Al termine di una sessione lavorativa, è bene:
\begin{enumerate}[1.]
\item eseguire l'aggiornamento del registro delle modifiche nella documentazione (se necessario)
\item committare e pushare, per evitare la perdita del lavoro effettuato
\item segnare le ore di lavoro svolte sulla
\href{https://docs.google.com/spreadsheets/d/12esX1ISWQOKM-fjuHTLmAzRN0cWltksn7eGiPsBHBI0/edit?usp=sharing}{Tabella delle ore lavorative} \cite{TabellaDelleOreLavorative}.
\item aggiornare lo stato di avanzamento su
\href{https://trello.com/b/hIEOGbE9/jawadruids}{Trello} \cite{TrelloJawaDruids}.
\end{enumerate}

\begin{thebibliography}{}
    \bibitem{GitFlowCheatseet}
    Daniel Kummer.
    \textit{Git flow cheatsheet - efficient branching using git-flow by Vincent Driessen}.\\
		\url{https://danielkummer.github.io/git-flow-cheatsheet/}.
		\bibitem{GitFlowWorkflowTutorial}
		Atlassian Bitbucket.
		\textit{Git flow workflow tutorial}.\\
		\url{https://www.atlassian.com/git/tutorials/comparing-workflows/gitflow-workflow}.
		\bibitem{GitCommitGuidelines}
		Cameron McKenzie, TechTarget.
		\textit{How to write a Git commit message properly with examples}.\\
		\url{https://www.theserverside.com/video/Follow-these-git-commit-message-guidelines}.
		\bibitem{TabellaDelleOreLavorative}
		Jawa Druids, 2020.
		\textit{Tabella delle Ore Lavorative}.\\
		\url{https://docs.google.com/spreadsheets/d/12esX1ISWQOKM-fjuHTLmAzRN0cWltksn7eGiPsBHBI0/edit?usp=sharing}.
		\bibitem{TrelloJawaDruids}
		Jawa Druids, 2020.
		\textit{Trello - dashboard Jawa Druids}.\\
		\url{https://trello.com/b/hIEOGbE9/jawadruids}.
\end{thebibliography}


\end{document}
